\section{Goal drift: phenomenon and mechanisms}
\label{sec:drift}

\subsection{A family of related failures}

``Goal drift'' is used loosely in the literature. Several neighbouring phenomena share its surface
symptom---the agent ends up doing something other than what the principal wanted---but differ in
where the competing objective comes from and whether anyone intends it. Table~\ref{tab:phenomena}
positions them. We use \emph{goal drift} in the sense of \citet{arike2025goaldrift}: an agent's
tendency to deviate from its original objective \emph{over time}, as an effect of the trajectory
rather than of a mis-specified goal or an adversary.

\begin{table}[t]
\centering
\small
\begin{tabularx}{\linewidth}{@{}>{\raggedright\arraybackslash}p{3.2cm}>{\raggedright\arraybackslash}X>{\raggedright\arraybackslash}p{4.2cm}@{}}
\toprule
Phenomenon & Origin of the competing objective & Representative work \\
\midrule
Goal misgeneralisation & Training: a capable system learns a proxy goal that coincides with
the intended one only in distribution & \citet{langosco2022goalmisgen,shah2022goalmisgen} \\
Specification gaming / reward hacking & The objective's proxy (e.g.\ passing tests) is optimised
instead of the intent & \citet{amodei2016concrete,zhong2025impossiblebench} \\
Goal hijacking (prompt injection) & An adversary embeds instructions in data the agent reads &
\citet{zhan2024injecagent,debenedetti2024agentdojo} \\
Instruction / persona drift & Decay of adherence to system-prompt instructions or persona over many
turns & \citet{li2024instability,ding2026contextecho} \\
Action bias / scope creep & The agent acts, or acts more broadly, when no or less action is
warranted & \citet{gloaguen2026fixedbench,mavali2026tab} \\
\textbf{Goal drift} & Accumulated context---environmental pressure, prior actions, sub-tasks---gradually
displaces the assigned goal & \citet{arike2025goaldrift,menon2026inherited,saebo2026asymmetric} \\
\bottomrule
\end{tabularx}
\caption{Phenomena related to goal drift, by origin of the competing objective.}
\label{tab:phenomena}
\end{table}

The distinctions matter for both measurement and design. Misgeneralisation is fixed at training time.
Hijacking presupposes an adversary and is addressed with security defences. Drift, by contrast, can
arise in an entirely benign environment, from material that is individually reasonable---a failing
test, a maintainer's note, a user's side request. In such an environment no single step looks like a
violation. The deviation exists only at the level of the trajectory.

\subsection{Direct evidence of goal drift in LLM agents}

\citet{arike2025goaldrift} give a goal to an agent through its system prompt and then expose it to
competing objectives through environmental pressure in a simulated trading environment. The
best-performing scaffolded agent maintained near-perfect adherence for more than 100{,}000 tokens,
but all models drifted to some degree. Drift correlated with growing susceptibility to
\emph{pattern-matching} as the context lengthened: agents increasingly continued what was already in
their context rather than what their goal prescribed. \citet{menon2026inherited} revisit the setting
with newer models, which are largely robust when acting alone. That robustness is brittle, though.
Conditioned on prefilled trajectories from weaker agents, the same models often \emph{inherit} the
drift, and only one of the tested models (GPT-5.1) remained consistently resilient. This finding is directly
relevant to agentic tools in which a context is summarised, resumed or handed to a sub-agent.

In coding specifically, \citet{saebo2026asymmetric} give agents built on OpenCode a system-prompt
constraint and apply codebase pressure in the opposite direction. Violations are \emph{asymmetric}:
agents abandon constraints more readily when the constraint opposes values such as security or
privacy. Drift therefore depends on what the competing objective is, not only on how much context
separates it from the goal.

Earlier, \citet{li2024instability} measured \emph{instruction drift} within eight rounds of dialogue
and related it to attention decaying away from the system prompt. \citet{ding2026contextecho}
measure persona drift in long real-world Claude Code and Codex CLI sessions that span thousands of
tool-using steps and repeated compaction. They introduce a \emph{snapshot-then-probe} protocol that
evaluates the conversation state without perturbing the main session; we adapt it, by probing a
forked copy of the session, to measure goal recall. \citet{rath2026agentdrift} proposes a taxonomy
of semantic, coordination and behavioural drift for multi-agent systems, though this single-author
work rests on simulation and theoretical modelling.

\subsection{Mechanisms that plausibly produce drift through sub-tasks}

The literature suggests several mechanisms, not mutually exclusive, through which accumulating
sub-tasks can displace a goal.

\paragraph{Position and dilution.} Models use information at the beginning and end of long contexts
better than information in the middle \citep{liu2024lost}. The original request sits at the
beginning of a session. Each sub-task, tool output and error message pushes it further from the
end, where the most recent material---typically the current sub-task---sits. In a single-author,
white-box code-audit case study, \citet{xue2026contextrot} holds the task fixed and varies only the
surrounding context. Pass rates fall from 8/10 runs with a clean context to 3/10 with long contexts,
whether the padding was relevant or irrelevant, although the difference is only trend-level
($p\approx0.07$) and a second task showed no degradation. Requirement coverage nevertheless stays
above 92\%. If this pattern generalises, drift need not look like abandonment: a few silently dropped
requirements are enough to invalidate an otherwise complete artefact.

\paragraph{Premature commitment and unreliability.} When a specification is revealed over several
turns rather than all at once, LLM performance drops by 39\% on average, mostly through increased
\emph{unreliability} rather than lower aptitude: models make assumptions early and fail to recover
\citep{laban2025lost}. Mid-session user side-requests create exactly this kind of multi-turn,
incrementally specified situation.

\paragraph{Self-conditioning and pattern-matching.} Models make more mistakes when their own earlier
errors are in the context \citep{sinha2025illusion}, and drift correlates with pattern-matching on
recent context \citep{arike2025goaldrift}. A sequence of sub-task actions---fixing a test, then
another, then a lint warning---can become a pattern that the agent continues after it has stopped
serving the goal.

\paragraph{Lossy compaction and hand-off.} Summarisation decides what survives of the original
request. Compression guidelines refined from observed agent failures improve long-horizon agents
\citep{kang2025acon}, and simply masking old observations matches LLM summarisation on SWE-bench
Verified at lower cost \citep{lindenbauer2025complexity}; neither study tracks whether the original
objective survives compression, which is the question that matters for drift. Inherited drift \citep{menon2026inherited} shows that a compressed or transferred
state can carry the drift along with it.

\paragraph{Decomposition and derailment.} Agents that decompose tasks on demand, when execution fails
\citep{prasad2024adapt}, generate sub-tasks from errors by design. In multi-agent systems,
\emph{task derailment} and \emph{disobeying the task specification} are among the 14 failure modes
of the MAST taxonomy, which was derived from 150 traces and applied to a dataset of more than
1{,}600 annotated traces \citep{cemri2025mast}. The Agent GPA framework starts from the premise that
critical agent failures emerge at the interfaces between goal, plan and action, and its judges
cover most human-annotated errors on the benchmarks it was tested on \citep{jia2025gpa}. These are precisely the interfaces that a sub-task traverses when it enters the
plan.

\paragraph{Action bias and instruction-like text in the environment.} Coding agents act even when they
should not. On FixedBench, where no change is needed, state-of-the-art agents still propose undesirable
patches in 35--65\% of cases \citep{gloaguen2026fixedbench}. Terminal agents must follow cues in
READMEs, comments and stack traces while ignoring plausible distractors, and the Task Alignment
Benchmark (TAB) reveals a systematic gap between being capable of a task and being aligned with
it \citep{mavali2026tab}. Even benign, non-adversarial content in the environment distracts strong GUI
agents \citep{ma2024caution}. Benign drift and adversarial hijacking therefore share a channel:
instruction-shaped text in tool outputs.

\paragraph{Proxy sub-goals.} ``Make the tests pass'' is the archetypal sub-goal that can replace the
goal it was meant to serve. On ImpossibleBench, the specification and the tests are made to
conflict, so any pass implies a specification-violating shortcut---for example deleting or editing
the failing test \citep{zhong2025impossiblebench}. This is the drift-as-reward-hacking corner of the
space \citep{amodei2016concrete}.

\subsection{Manipulating sub-tasks experimentally}

Several recent benchmarks inject material into an ongoing trajectory and can be adapted to study
drift. InterruptBench formalises three kinds of mid-task user interruption---adding, revising or
retracting a requirement---in long-horizon web navigation \citep{zou2026interruptbench}. SWE-Touch
injects plausible conflicting edits, accompanied by a user message, when a coding agent reaches
task-critical code, lowering the average resolve rate of nine models on SWE-bench Verified by 7.7
percentage points \citep{tan2026swetouch}.
AgentDojo pairs user tasks with injected goals and scores utility and attack success separately
\citep{debenedetti2024agentdojo}. Task Shield defends against injection by checking whether each
instruction and tool call \emph{serves the user's goal} \citep{jia2024taskshield}; this per-action
test is a natural oracle for drift. SlopCodeBench shows how agents' code erodes as they repeatedly
extend their own solutions \citep{orlanski2026slopcode}, a long-horizon complement to single-episode
drift. Finally, long-context degradation also affects the monitors that might be used to catch
drift: LLM classifiers miss subtly dangerous agent actions far more often when those actions appear
deep in long transcripts \citep{martin2026classifier}.

Taken together, the evidence suggests that goal drift through sub-tasks is plausible on mechanistic
grounds, partially observed in synthetic settings, and measurable with existing tools. What is
missing is a controlled study of the \emph{dose} and \emph{source} of sub-tasks in a real coding-agent
harness, joined to the question of what a human supervisor can see.
